% =========================================================
% SECCIÓN 2.1 - INTRODUCCIÓN
% =========================================================

La bioconversión de residuos orgánicos mediante \textit{Hermetia illucens} opera, en la mayoría de las implementaciones industriales y académicas, como un proceso de ``caja negra'': se controlan los insumos (sustrato, densidad larval, humedad) y se cuantifican los productos (biomasa, frass), pero los flujos gaseosos intermedios permanecen sin caracterizar dinámicamente \parencite{jenkinsProcessingPoultryManure2025, fuhrmannComprehensiveIndustryrelevantBlack2025}. Esta opacidad metrológica impide establecer factores de emisión específicos, validar balances de carbono en tiempo real y detectar condiciones operativas subóptimas ---como la anaerobiosis transitoria--- que degradan tanto el desempeño ambiental como la eficiencia zootécnica del sistema \parencite{bekkerImpactSubstrateMoisture2021, rossiEstimatingDynamicsGreenhouse2024}.

Frente a esta brecha, la instrumentación basada en sensores de infrarrojo no dispersivo (NDIR) y arquitecturas IoT de bajo costo ha emergido como una alternativa viable para el monitoreo continuo de \ce{CO2} y \ce{CH4} en entornos agroindustriales \parencite{rajRealTimeEstimation2023, duobieneDevelopmentWirelessSensor2022, yavariArtEMonArtificialIntelligence2023}. Sin embargo, la transferencia directa de hardware diseñado para ambientes controlados (invernaderos, almacenes) a reactores de bioconversión expone limitaciones críticas: deriva instrumental por condensación, saturación de señal en presencia de amoniaco y compuestos orgánicos volátiles, y una incertidumbre que frecuentemente excede los umbrales exigidos por protocolos de Medición, Reporte y Verificación (MRV) \parencite{ahmedAgricultureClimateChange2020, herrinCollateralDataQuality2021}. Por tanto, la viabilidad de estos sensores no depende únicamente de sus especificaciones de fabricante, sino de una caracterización metrológica local que incluya curvas de calibración frente a patrones certificados, compensación cruzada por variables ambientales y una estimación rigurosa de la incertidumbre combinada \parencite{vanIntegrationInternetofThingsSustainable2022}.

En este contexto, el presente capítulo aborda el diseño experimental y la caracterización del módulo de adquisición que alimentará el sistema híbrido descrito en los capítulos posteriores. Se estructura en cinco bloques. Primero, se describe el reactor batch y las condiciones de laboratorio en las que se ejecutaron los ciclos de bioconversión (\autoref{sec:cap2_sistema_experimental}). Segundo, se detallan los sensores NDIR seleccionados, las variables ambientales auxiliares y la arquitectura del nodo IoT con computación en el borde (\autoref{sec:cap2_sensores}). Tercero, se presenta el protocolo de medición en ventanas cerradas ---estandarizado para estimar tasas de emisión en sistemas de residuos sólidos--- incluyendo la conversión de concentración (ppm) a flujo másico (\si{\gram\per\minute}) (\autoref{sec:cap2_protocolo}). Cuarto, se exponen las curvas de calibración, los modelos de compensación cruzada y el presupuesto de incertidumbre metrológica (\autoref{sec:cap2_calibracion}). Finalmente, se sintetizan los resultados de la caracterización instrumental y se establece la trazabilidad del dato crudo que servirá de entrada al modelo bioenergético del Capítulo 3 y al estimador híbrido del Capítulo 4.

El objetivo metrológico central es garantizar que la señal de los sensores NDIR presente una correlación de Pearson $R > 0{,}85$ y un error porcentual absoluto medio (MAPE) inferior al \SI{15}{\percent} respecto a un instrumento de referencia, bajo las condiciones de temperatura y humedad relativas propias del cultivo de \textit{H. illucens} (\SI{30}{\celsius} $\pm$\,\SI{2}{\celsius}, HR \SI{60}{\percent}--\SI{80}{\percent}). El cumplimiento de este umbral constituye la hipótesis H1 y habilita la integración posterior con el modelo matemático sin propagación excesiva de errores instrumentales.
